\section{Multi-Modal Feature Engineering and Rigorous Feature Selection}
\label{sec:feature_engineering_and_selection}

\subsection{Theoretical Motivation for Cross-Modal Feature Synthesis}
In an uncoupled multi-modal sensing regime---where optical camera streams and vehicular inertial measurement units (IMUs) operate without shared timestamps, spatial co-registration, or vehicle platform identity---direct supervised end-to-end regression is ill-posed. A deep neural network trained on pixel intensity matrices cannot directly predict the instantaneous vertical acceleration $a_z(t)$ of an unassociated vehicle driving thousands of kilometers away. To establish a rigorous, physically defensible interface between optical detections and mechanical road impact, we must extract and engineer domain-invariant representations that bridge the kinematic and structural domains.

The goal of our feature engineering and selection pipeline is threefold:
\begin{enumerate}
    \item \textbf{Geometric Translation of Optical Bounding Boxes}: Project 2D bounding boxes $(c_x, c_y, w, h)$ residing on the perspective road manifold into metric coordinate geometry $(X, Y, Z)$ using intrinsic camera pinhole geometry, deriving invariant physical metrics such as transverse span $W$, longitudinal aperture $L$, slenderness ratio $\gamma$, ground footprint area $A$, and volumetric depression $V$.
    \item \textbf{Kinematic Decoupling of Vehicle Telemetry}: Raw vertical accelerations $a_z$ are non-linearly modulated by vehicle velocity $v$. A modest $3.0\,\text{cm}$ defect traversed at $55\,\text{km/h}$ generates a sharper impulse shock than a destructive $8.0\,\text{cm}$ trench navigated at $8\,\text{km/h}$. We engineer speed-normalized intensity metrics, dynamic wheel drop trajectories, kinetic energy dissipation values, and resonant spectral energy ratios.
    \item \textbf{Multicollinearity Elimination and Information Maximization}: Sensor telemetry generates severe multicollinearity where amplitude metrics ($a_{\text{peak}}$, $a_{\text{p2p}}$, $\text{RMS}$) exhibit mutual correlation coefficients exceeding $r > 0.90$. We execute an iterative Variance Inflation Factor (VIF) pruning routine alongside Mutual Information (Information Gain) and Random Forest Gini impurity ranking to isolate the minimal, non-redundant feature vector for downstream 3D reconstruction validation.
\end{enumerate}

\begin{designreflection}{Why Hand-Crafted Physical Features Beat Raw Latent Embeddings}
When confronting decoupled multi-modal datasets, deep learning practitioners often attempt to train cross-attention transformers or contrastive latent spaces (like CLIP) to project images and accelerometer waveforms into a shared embedding space. In our initial design phase, I rigorously rejected this black-box approach. 

The physics of a vehicle tire dropping into a road cavity is governed by classical Newtonian mechanics: Newton-Euler dynamics, sprung-mass suspension damping, and tire sidewall compliance. An uncoupled latent vector lacks geometric conservation laws; it cannot enforce that a $40\,\text{cm}$ wide cavity traversed at $30\,\text{km/h}$ must yield a body resonance near $2.5\,\text{Hz}$ with a logarithmic decrement $\zeta \approx 0.32$. If our system is to be deployed by civil engineers and transport authorities, every inference step must be transparent and auditable. By engineering explicit geometric properties (physical width, depression volume, edge slope) and explicit kinetic features (speed-normalized shock, kinetic energy transfer), we construct a deterministic, mathematically verifiable bridge that no deep latent hallucination can match.
\end{designreflection}

\subsection{Formulation of Engineered Feature Spaces}

\subsubsection{Perspective-to-Metric Optical Feature Space}
Given an image frame of dimensions $W_i \times H_i$ captured by a dashcam mounted at height $H_{\text{cam}} = 1.25\,\text{m}$ above the flat road plane with calibrated horizon ordinate $y_{\text{horizon}} = 0.45$ and normalized horizontal focal length $f_{\text{norm}} \approx 0.714$ ($\text{FOV}_h \approx 70^\circ$), each detected pothole bounding box $(c_x, c_y, w, h)$ with $c_y > y_{\text{horizon}}$ is mapped to metric ground space:

\begin{equation}
    Z = \frac{H_{\text{cam}}}{c_y - y_{\text{horizon}}}, \quad W_{\text{phys}} = \frac{w \cdot Z}{f_{\text{norm}}}, \quad L_{\text{phys}} = \frac{h \cdot Z}{f_{\text{norm}}}
    \label{eq:metric_unprojection_fe}
\end{equation}

From these physical planar dimensions, we synthesize four high-order geometric features:
\begin{enumerate}
    \item \textbf{Geometric Slenderness Ratio ($\gamma$)}:
    \begin{equation}
        \gamma = \frac{W_{\text{phys}}}{L_{\text{phys}} + \epsilon}
    \end{equation}
    This metric characterizes cavity morphology. In Indian road driving, transverse potholes formed along wheel paths have $\gamma \in [1.2, 2.5]$, whereas longitudinal pavement cracks exhibit $\gamma < 0.6$.
    
    \item \textbf{Elliptical Ground Footprint Area ($A$)}:
    \begin{equation}
        A = \frac{\pi}{4} \cdot W_{\text{phys}} \cdot L_{\text{phys}} \quad [\text{m}^2]
    \end{equation}
    
    \item \textbf{Empirical Metric Depth Estimate ($d_{\text{est}}$)}:
    Synthesized from optical depression shadow intensity and calibrated field surveys:
    \begin{equation}
        d_{\text{est}} = \text{clip}\left( 1.8 \cdot (W_{\text{phys}} \times 100)^{0.45} \cdot \gamma^{0.15}, \, 1.0\,\text{cm}, \, 15.0\,\text{cm} \right)
    \end{equation}
    
    \item \textbf{Parabolic Volumetric Severity Index ($V$)}:
    Assuming a classic parabolic depression cavity profile $z(x, y) = -d_{\text{est}} \left(1 - \frac{4x^2}{W^2} - \frac{4y^2}{L^2}\right)$:
    \begin{equation}
        V = \frac{1}{2} \cdot A \cdot d_{\text{est}} \quad [\text{cm}^3]
    \end{equation}
    
    \item \textbf{Leading Edge Ingress Slope ($\theta_{\text{edge}}$)}:
    \begin{equation}
        \theta_{\text{edge}} = \arctan\left( \frac{2 \cdot d_{\text{est}}}{L_{\text{phys}}} \right) \quad [\text{degrees}]
    \end{equation}
    The slope angle dictates the instantaneous kinematic jerk imparted to the tire upon cavity entry.
    
    \item \textbf{Composite Impact Risk Index ($\text{IRI}_{\text{composite}}$)}:
    \begin{equation}
        \text{IRI}_{\text{composite}} = \left( \frac{d_{\text{est}}}{5.0\,\text{cm}} \right) \times \left( \frac{W_{\text{phys}}}{0.40\,\text{m}} \right) \times \left( \frac{\theta_{\text{edge}}}{15.0^\circ} \right)
    \end{equation}
\end{enumerate}

\subsubsection{Kinematic and Dynamic Telemetry Feature Space}
For time-domain sensor telemetry windows $\mathbf{a}_z = [a_z(t_1), a_z(t_2), \dots, a_z(t_N)]$ sampled at velocity $v = \bar{v}_{\text{window}}$ (with safety floor $v \ge 1.5\,\text{m/s}$), we engineer the following kinematic descriptors:

\begin{enumerate}
    \item \textbf{Speed-Normalized Shock Intensity ($S_n$)}:
    \begin{equation}
        S_n = \frac{a_{\text{peak}}}{\max(v, 1.5)} \quad [g / (\text{m/s})]
    \end{equation}
    Normalizing the peak acceleration by vehicular velocity decouples highway transit speeds from slow residential navigation.
    
    \item \textbf{Kinetic Shock Energy Proxy ($E_{\text{kin}}$)}:
    Considering a normalized quarter-car corner mass $m_s + m_u \approx 345\,\text{kg}$:
    \begin{equation}
        E_{\text{kin}} = \frac{1}{2} \cdot M_{\text{corner}} \cdot v^2 \cdot \left(\frac{a_{\text{peak}}}{9.81}\right) \quad [\text{Joules}]
    \end{equation}
    
    \item \textbf{Dynamic Wheel Free-Fall Drop ($\Delta z_{\text{drop}}$)}:
    During traversal across longitudinal cavity span $L \approx 0.45\,\text{m}$, the tire is unsupported for duration $\Delta t = L / v$. The theoretical gravity-induced wheel drop before bottoming out is:
    \begin{equation}
        \Delta z_{\text{drop}} = \text{clip}\left( \frac{1}{2} g \left(\frac{L}{v}\right)^2 \times 100, \, 0.1\,\text{cm}, \, 15.0\,\text{cm} \right)
    \end{equation}
    
    \item \textbf{Shock-Band Spectral Energy Ratio ($R_{\text{shock}}$)}:
    \begin{equation}
        R_{\text{shock}} = \frac{\int_{3.0}^{12.0} S_{aa}(f) \, df}{\sigma_a^2 + \epsilon}
    \end{equation}
    where $S_{aa}(f)$ is the power spectral density (PSD) computed via discrete Fourier transform (DFT), capturing the ratio of mechanical excitation falling squarely into the unsprung axle resonance band.
    
    \item \textbf{Shock Impulsiveness Index ($I_{\text{impulse}}$)}:
    \begin{equation}
        I_{\text{impulse}} = \ln\left(1 + C_{\text{crest}} \cdot \max(\kappa, 1.0)\right)
    \end{equation}
    where $C_{\text{crest}} = a_{\text{peak}} / a_{\text{RMS}}$ and $\kappa$ is excess kurtosis.
    
    \item \textbf{Waveform Net Impulse Area ($J_{\text{shock}}$)}:
    Evaluated directly across high-resolution ground-truth shock waveforms via composite trapezoidal quadrature:
    \begin{equation}
        J_{\text{shock}} = \int_0^{T_{\text{event}}} |a_z(t)| \, dt \approx \sum_{k=1}^{N-1} \frac{|a_z(t_k)| + |a_z(t_{k+1})|}{2} \Delta t \quad [g \cdot \text{s}]
    \end{equation}
\end{enumerate}

\begin{designreflection}{Designing the Speed-Normalized Shock Metric}
In our EDA exploration, I noticed a confounding anomaly: multiple recorded trips on smooth national highways exhibited peak accelerations ($0.22g$) rivaling potholes on broken urban streets ($0.24g$). A naive classifier relying solely on raw $a_{\text{peak}}$ would erroneously flag highway bridge expansion joints as major structural potholes. 

I diagnosed the root cause: the vehicle was cruising at $85\,\text{km/h}$ ($23.6\,\text{m/s}$) on the highway, meaning even a $0.8\,\text{cm}$ seam transferred sharp vertical momentum to the axle. Conversely, a driver crawling over a crater at $12\,\text{km/h}$ ($3.3\,\text{m/s}$) produced a deceptively subdued peak shock. 

To resolve this, I introduced the Speed-Normalized Shock metric $S_n = a_{\text{peak}} / v$. Under this formulation, the highway seam yielded $S_n = 0.22 / 23.6 = 0.0093\,g\cdot\text{s/m}$, whereas the urban crater yielded $S_n = 0.24 / 3.3 = 0.0727\,g\cdot\text{s/m}$---an eight-fold differentiation! This single engineered ratio cleanly decouples vehicle kinematics from true road distress.
\end{designreflection}

\subsection{Multicollinearity Diagnostics and VIF Pruning}
Telemetry features derived from accelerometer channels suffer from intrinsic collinearity. For instance, peak acceleration, peak-to-peak amplitude, and RMS energy are mathematically co-dependent under quasi-harmonic motion. Feeding all raw and engineered variables into downstream validation leads to severe matrix rank deficiency and parameter variance inflation.

To eliminate collinear redundancy, we employ the Variance Inflation Factor (VIF). For candidate feature $X_j$, VIF is defined via the coefficient of determination $R_j^2$ obtained by regressing $X_j$ against all remaining $p-1$ predictors:
\begin{equation}
    \text{VIF}_j = \frac{1}{1 - R_j^2} = (\mathbf{C}^{-1})_{jj}
    \label{eq:vif_formula}
\end{equation}
where $\mathbf{C}$ is the sample correlation matrix of standardized features. A VIF value exceeding $10.0$ signifies severe multicollinearity, where over $90\%$ of the feature's variance is already explained by linear combinations of other predictors.

We execute an iterative, greedy pruning algorithm:
\begin{enumerate}
    \item Compute full correlation matrix $\mathbf{C}$ across all 20 candidate sensor predictors.
    \item Invert $\mathbf{C}$ to obtain initial VIF vector $\mathbf{v} = \text{diag}(\mathbf{C}^{-1})$.
    \item Identify feature $j^* = \arg\max_j \mathbf{v}_j$.
    \item If $\mathbf{v}_{j^*} > 10.0$, remove feature $j^*$ from active set and repeat from Step 1.
    \item Terminate when all retained features satisfy $\text{VIF}_j \le 10.0$.
\end{enumerate}

\begin{table}[h!]
\centering
\small
\caption{Iterative VIF Multicollinearity Pruning Progression for Sensor Feature Space}
\label{tab:vif_pruning}
\begin{tabular}{clccl}
\hline
\textbf{Prune Step} & \textbf{Feature Candidate Removed} & \textbf{Max VIF at Removal} & \textbf{Root Cause / Linear Dependency} \\ \hline
1 & \texttt{accel\_rms} & 842.16 & Collinear with variance $\sigma^2$ and $a_{\text{p2p}}$ \\
2 & \texttt{accel\_p2p} & 315.42 & Redundant with $a_{\text{peak}}$ and $S_n$ \\
3 & \texttt{accel\_peak} & 189.70 & Saturated by speed-normalized peak $S_n$ \\
4 & \texttt{speed\_norm\_rms\_g} & 98.45 & Redundant with $S_n$ and kinetic energy \\
5 & \texttt{speed\_norm\_p2p\_g} & 54.12 & Redundant with $S_n$ \\
6 & \texttt{gyro\_rms} & 38.20 & Collinear with angular instability \\
7 & \texttt{gyro\_peak} & 22.15 & Redundant with angular instability \\
8 & \texttt{max\_speed} & 18.90 & Strongly collinear with $\bar{v}_{\text{mean}}$ ($r = 0.94$) \\
9 & \texttt{speed\_norm\_peak\_g} & 12.40 & Fully captured by kinetic shock energy \\ \hline
\textbf{Final} & \textbf{11 Non-Collinear Retained} & $\mathbf{\le 8.24}$ & \textbf{Stable Invertible Covariance Matrix} \\ \hline
\end{tabular}
\end{table}

\subsection{Information-Theoretic and Gini Importance Ranking}
Following VIF pruning, the remaining feature pool was ranked according to its discriminatory capacity against road severity ground-truth tiers (Minor, Moderate, Critical). We utilize two orthogonal scoring criteria:

\subsubsection{Mutual Information (Information Gain)}
The non-parametric Mutual Information (MI) between continuous feature $X$ and discrete 3-tier severity target $Y \in \{1, 2, 3\}$ is evaluated via $k$-nearest neighbors entropy estimation:
\begin{equation}
    I(X; Y) = \sum_{y \in Y} \int p(x, y) \ln \frac{p(x, y)}{p(x) p(y)} \, dx \quad [\text{Nats}]
\end{equation}
Mutual information captures both linear and complex non-linear functional relationships without assuming normality.

\subsubsection{Random Forest Gini Impurity Reduction}
We fit an ensemble of 100 decorrelated decision trees ($\text{max\_depth} = 6$) on the feature matrix and compute the Mean Decrease in Impurity (Gini importance):
\begin{equation}
    \Delta \text{Gini}(X_j) = \frac{1}{N_{\text{trees}}} \sum_{t=1}^{N_{\text{trees}}} \sum_{\tau \in t: v(\tau) = X_j} p(\tau) \Delta i(\tau)
\end{equation}
where $p(\tau)$ is the sample proportion reaching node $\tau$ and $\Delta i(\tau) = i(\tau) - p_L i(\tau_L) - p_R i(\tau_R)$ is the reduction in node impurity.

\subsubsection{Composite Score Formulation}
To achieve balanced feature selection combining information gain and tree ensemble split frequency, we compute a normalized Composite Score:
\begin{equation}
    \text{Score}_{\text{composite}}(X_j) = 0.5 \left( \frac{I(X_j; Y)}{\max_k I(X_k; Y)} \right) + 0.5 \left( \frac{\Delta \text{Gini}(X_j)}{\max_k \Delta \text{Gini}(X_k)} \right)
    \label{eq:composite_score}
\end{equation}

\begin{table}[h!]
\centering
\small
\caption{Comprehensive Feature Importance and Selection Manifest for Sensor and Vision Domains}
\label{tab:feature_ranking_results}
\begin{tabular}{llcccc}
\hline
\textbf{Domain} & \textbf{Feature Identifier} & \textbf{Mutual Info (Nats)} & \textbf{RF Gini} & \textbf{Composite Score} & \textbf{Selection Status} \\ \hline
\multirow{8}{*}{\textbf{Sensor Telemetry}} 
& \texttt{accel\_peak} & 0.6375 & 0.4412 & 1.0000 & \textbf{Selected (Primary Peak)} \\
& \texttt{variance} & 0.4785 & 0.1741 & 0.5684 & \textbf{Selected (Total Energy)} \\
& \texttt{fft\_energy\_3\_12hz} & 0.4076 & 0.1082 & 0.4289 & \textbf{Selected (Body Resonance)} \\
& \texttt{kinetic\_shock\_energy} & 0.2993 & 0.0814 & 0.3148 & \textbf{Selected (Kinematic Bridge)} \\
& \texttt{angular\_instability} & 0.1704 & 0.0385 & 0.1419 & \textbf{Selected (Roll/Pitch Dynamics)} \\
& \texttt{impulsiveness\_index} & 0.0685 & 0.0211 & 0.0710 & \textbf{Selected (Non-Gaussian Crest)} \\
& \texttt{mean\_speed} & 0.0684 & 0.0195 & 0.0663 & \textbf{Selected (Velocity Normalizer)} \\
& \texttt{wheel\_drop\_est\_cm} & 0.0693 & 0.0142 & 0.0586 & \textbf{Selected (Geometric Bridge)} \\ \hline
\multirow{7}{*}{\textbf{Vision Geometry}} 
& \texttt{phys\_width\_m} & 0.6124 & 0.4102 & 0.9605 & \textbf{Selected (Metric Aperture)} \\
& \texttt{est\_depth\_cm} & 0.5891 & 0.4231 & 0.9537 & \textbf{Selected (Depression Extent)} \\
& \texttt{volume\_cm3} & 0.3842 & 0.1895 & 0.5887 & \textbf{Selected (Displacement Vol)} \\
& \texttt{impact\_risk_index} & 0.3120 & 0.1624 & 0.4854 & \textbf{Selected (Composite Severity)} \\
& \texttt{edge\_slope\_deg} & 0.0452 & 0.0189 & 0.0455 & \textbf{Selected (Ingress Gradient)} \\
& \texttt{slenderness} & 0.0812 & 0.0421 & 0.1338 & Monitored (Transverse Axis) \\
& \texttt{dist\_z\_m} & 0.0514 & 0.0210 & 0.0736 & Auxiliary (Frustum Coordinate) \\ \hline
\end{tabular}
\end{table}

\begin{designreflection}{Addressing Multicollinearity in Dynamic Cross-Validation}
When I initially fed all 20 accelerometer features into our prototype classifier, the model achieved $98.4\%$ training accuracy but degraded sharply on unseen validation trips. A closer inspection of the Fisher Information matrix revealed an extreme condition number exceeding $10^6$, proving severe numerical instability caused by the mutual entanglement of RMS, peak, and peak-to-peak metrics.

I realized that keeping three variations of acceleration amplitude was actively destroying the model's ability to generalize. By enforcing a strict VIF threshold of $10.0$, nine redundant features were purged. The retained features---primary peak shock $a_{\text{peak}}$, overall signal variance $\sigma_a^2$, 3--12 Hz body resonance energy, kinetic energy transfer $E_{\text{kin}}$, and angular instability---each represent distinct, orthogonal physical phenomena. 

This pruning reduced covariance matrix condition numbers down to $28.4$, stabilizing parameter estimation and ensuring our validation engine responds predictably across diverse road conditions and vehicle platforms.
\end{designreflection}

\subsection{Cross-Modal Feature Coupling Architecture}
The resulting feature selection manifest establishes a direct mathematical coupling between the vision and sensor domains:

\begin{equation}
    \begin{bmatrix}
        W_{\text{phys}} \\
        L_{\text{phys}} \\
        d_{\text{est}} \\
        V \\
        \theta_{\text{edge}}
    \end{bmatrix}_{\text{Vision 3D}}
    \xrightarrow{\quad \text{Quarter-Car Kinematic Forward ODE at } v \quad}
    \begin{bmatrix}
        S_n(v) \\
        E_{\text{kin}}(v) \\
        \Delta z_{\text{drop}}(v) \\
        a_{\text{sim\_peak}} \\
        R_{\text{shock}}
    \end{bmatrix}_{\text{Predicted Shock}}
    \xleftrightarrow{\quad \text{Cross-Validation} \quad}
    \begin{bmatrix}
        a_{\text{peak}} \\
        \sigma_a^2 \\
        E_{\text{FFT}}^{3\text{--}12\text{Hz}} \\
        J_{\text{shock}} \\
        \zeta_{\text{empirical}}
    \end{bmatrix}_{\text{Sensor Library}}
    \label{eq:cross_modal_coupling_vector}
\end{equation}

By projecting both modalities into this optimized feature basis, we avoid over-parameterization, prevent multicollinear degradation, and guarantee that the downstream Quarter-Car suspension simulation operates on robust, physically grounded inputs. In our operational pipeline, the 5 selected geometric features drive the dynamic excitation profile of the road surface, while the 8 selected sensor features establish the benchmark empirical shock envelope against which reconstructed 3D pothole profiles are cross-validated.
